% =============================================================================
%                    WEEK 11: LLM EVALUATION AND SAFETY
% =============================================================================
\section{Week 11: LLM Evaluation and Safety}

% -----------------------------------------------------------------------------
%                              11.1 TOPIC SUMMARY
% -----------------------------------------------------------------------------
\subsection{Topic Summary}

\begin{keybox}[Learning Objectives]
By the end of this week, you will be able to:
\begin{enumerate}
  \item Learn about common \textbf{benchmarks} used to evaluate LLMs.
  \item Explore different \textbf{methods} to evaluate LLMs (automatic, human, LLM-as-judge).
  \item Understand key \textbf{metrics and challenges} in LLM evaluation.
  \item Understand key \textbf{safety risks} in LLMs (bias, toxicity, jailbreaking).
  \item Evaluate \textbf{defense strategies} for improving LLM safety.
  \item Critically assess \textbf{trade-offs and limitations} of current safety approaches.
\end{enumerate}
\end{keybox}

\separator

\subsubsection*{The Big Picture}

This week covers two critical aspects of LLMs: (1) \textbf{Evaluation}---how to measure LLM capabilities across tasks, benchmarks, and metrics; and (2) \textbf{Safety}---understanding vulnerabilities (bias, toxicity, jailbreaking) and defense strategies. Both are essential for responsible LLM deployment.

\separator

\subsubsection*{What Examiners Like to Ask}

\begin{itemize}
  \item Why might high benchmark scores not reflect true reasoning ability?
  \item Explain \textbf{Goodhart's Law} in the context of LLM leaderboards.
  \item Compare \textbf{text overlap metrics} (BLEU, ROUGE) vs.\ \textbf{semantic similarity} (BERTScore).
  \item Explain \textbf{LLM-as-a-Judge} and its challenges.
  \item What does performance drop on \textbf{GSM-Symbolic} suggest about LLM reasoning?
  \item Define \textbf{representational} vs.\ \textbf{allocational} harm.
  \item Describe bias evaluation methods: embedding-based, probability-based, generation-based.
  \item Compare \textbf{black-box} vs.\ \textbf{white-box} jailbreaking techniques.
  \item Explain defense strategies: SmoothLLM, GradSafe.
\end{itemize}

\bigseparator

% -----------------------------------------------------------------------------
%                 11.2 OVERVIEW + CORE CONCEPTS + EXTENDED THEORY
% -----------------------------------------------------------------------------
\subsection{Overview + Core Concepts + Extended Theory}

\subsubsection*{Roadmap}
\begin{enumerate}
  \item Why evaluation matters
  \item Evaluation tasks and benchmarks
  \item Evaluation methods (automatic, human, LLM-as-judge)
  \item Evaluation metrics
  \item Evaluation challenges
  \item Safety overview
  \item Social bias in LLMs
  \item Toxicity in LLMs
  \item Jailbreaking attacks (black-box and white-box)
  \item Defense strategies
\end{enumerate}

\bigseparator

% =============================================================================
%                              PART 1: EVALUATION
% =============================================================================

\subsubsection{Why Evaluation Matters}

\begin{defbox}[Importance of LLM Evaluation]
\begin{itemize}
  \item Understand \textbf{true capabilities} beyond raw scores
  \item Ensure \textbf{fairness, robustness, and safety}
  \item Avoid \textbf{overfitting to benchmarks}
  \item Enable meaningful \textbf{model comparison}
\end{itemize}
\end{defbox}

\begin{tipbox}[Benchmark Saturation]
When models reach $>90\%$ on benchmarks like AIME, the benchmark no longer distinguishes real progress. High scores don't necessarily mean human-level reasoning.
\end{tipbox}

\bigseparator

% --- Evaluation Tasks ---
\subsubsection{Evaluation Tasks}

\begin{center}
\renewcommand{\arraystretch}{1.3}
\begin{tabular}{@{} l p{7cm} @{}}
\toprule
\textbf{Category} & \textbf{Examples} \\
\midrule
Language Understanding & Reading comprehension, NLI, summarisation, sentiment \\
Knowledge \& Reasoning & MMLU, BIG-bench, mathematical reasoning, fact verification \\
Dialogue \& Interaction & Instruction following, helpfulness, harmlessness (HHH) \\
Safety \& Robustness & Toxicity detection, bias testing, adversarial robustness \\
Multimodal & Image+text reasoning, visual QA, chart/table reasoning \\
Specialised & Theory of Mind, emotion understanding, tool use \\
\bottomrule
\end{tabular}
\end{center}

\bigseparator

% --- Evaluation Benchmarks ---
\subsubsection{Evaluation Benchmarks}

\begin{defbox}[Benchmark Types]
\textbf{General benchmarks:} MMLU, BIG-bench, HELM, MT-Bench, AGIEval, C-Eval

\textbf{Domain-specific:} MultiMedQA (medical), CUAD (legal), ChemBench (science), SafetyBench

\textbf{Multimodal:} MMBench, VQAv2, MathVista, ToolBench
\end{defbox}

\begin{exbox}[Key Benchmarks]
\begin{itemize}
  \item \textbf{MMLU:} 57 subjects, 15,908 MCQs covering multitask knowledge
  \item \textbf{ARC-AGI:} Abstract reasoning requiring semantic interpretation of symbols
  \item \textbf{Chatbot Arena:} Blind pairwise comparisons where users vote on anonymous LLM responses
  \item \textbf{GSM-Symbolic:} Tests whether LLMs rely on pattern matching vs.\ true reasoning
\end{itemize}
\end{exbox}

\bigseparator

% --- Evaluation Methods ---
\subsubsection{Evaluation Methods}

\begin{thmbox}[Three Evaluation Approaches]
\textbf{1. Automatic Evaluation:}
\begin{itemize}
  \item Accuracy metrics: Exact Match, F1, BLEU, ROUGE
  \item Calibration, fairness, robustness metrics
\end{itemize}

\textbf{2. Human Evaluation:}
\begin{itemize}
  \item Expert assessment on accuracy, relevance, helpfulness
  \item Crowdsourced evaluation from multiple non-experts
  \item Comparative/pairwise evaluation
\end{itemize}

\textbf{3. LLM-as-a-Judge:}
\begin{itemize}
  \item Single model judging: Strong LLM evaluates other outputs
  \item Multi-model consensus: Multiple LLMs as judges, aggregate scores
  \item Constitutional AI evaluation: Models trained specifically for evaluation
\end{itemize}
\end{thmbox}

\bigseparator

% --- Evaluation Metrics ---
\subsubsection{Evaluation Metrics}

\begin{defbox}[Text Overlap Metrics]
\begin{itemize}
  \item \textbf{Exact Match (EM):} \% of answers exactly matching ground truth
  \item \textbf{BLEU:} Precision-focused n-gram overlap
  \item \textbf{ROUGE:} Recall-focused n-gram overlap (ROUGE-N, ROUGE-L)
  \item \textbf{METEOR:} Handles synonyms and word-order variations
\end{itemize}

\textbf{Problem:} These ignore semantic similarity between reference and candidate.
\end{defbox}

\begin{defbox}[Semantic Similarity Metrics]
\begin{itemize}
  \item \textbf{BERTScore:} Compares token embeddings from BERT
  \item \textbf{SentenceBERT:} Encodes full sentences, measures cosine similarity
  \item \textbf{BLEURT:} Trained to predict human evaluation scores
\end{itemize}
\end{defbox}

\begin{exbox}[Metric Comparison]
\textbf{Candidate:} ``The cat sits on the mat.''

\textbf{Reference:} ``The cat is sitting on the mat.''

\begin{center}
\begin{tabular}{@{} l c @{}}
\toprule
\textbf{Metric} & \textbf{Score} \\
\midrule
Exact Match & 0 \\
Token F1 & 0.77 \\
BLEU & 0.42 \\
ROUGE-1 & 0.77 \\
METEOR & 0.88 \\
\bottomrule
\end{tabular}
\end{center}

METEOR captures synonyms (``sits'' $\approx$ ``sitting'') better than BLEU.
\end{exbox}

\bigseparator

% --- LLM-as-a-Judge ---
\subsubsection{LLM-as-a-Judge}

\begin{defbox}[LLM-as-a-Judge Approaches]
\begin{itemize}
  \item \textbf{Pointwise:} Score each response independently
  \item \textbf{Pairwise:} Compare two responses, choose better one
  \item \textbf{Listwise:} Rank multiple responses
  \item \textbf{Multi-model debate:} Multiple LLMs argue to reach consensus
\end{itemize}
\end{defbox}

\begin{tipbox}[LLM-as-Judge Challenges]
\begin{itemize}
  \item \textbf{Bias:} Prefers longer, well-formatted, authoritative responses
  \item \textbf{Egocentric bias:} Prefers outputs similar to its own style
  \item \textbf{Adversarial vulnerability:} Susceptible to prompt injection
  \item \textbf{Scaling:} Longer reasoning improves accuracy but increases compute
\end{itemize}
\end{tipbox}

\bigseparator

% --- Evaluation Challenges ---
\subsubsection{Evaluation Challenges}

\begin{defbox}[GSM vs.\ GSM-Symbolic]
\textbf{Finding:} LLMs perform well on GSM-8K but poorly on GSM-Symbolic (same problems with different numbers/names).

\textbf{Implication:} Models may rely on \textbf{pattern matching} rather than true reasoning---they memorise solutions rather than understand the logic.
\end{defbox}

\begin{defbox}[Clause Order Sensitivity]
\textbf{Finding:} Reordering premises in logical problems causes significant performance drops.

\textbf{Implication:} LLMs rely on \textbf{word order patterns} rather than deep logical structure. Reordering shouldn't affect logical reasoning, but it does.
\end{defbox}

\begin{defbox}[Goodhart's Law]
\textbf{Statement:} ``When a measure becomes a target, it ceases to be a good measure.''

\textbf{In LLM context:} When models optimise for leaderboard rankings, they may overfit to benchmark tasks rather than improve general reasoning.

\textbf{Risk:} Developers optimise for scores instead of real-world quality.
\end{defbox}

\bigseparator

% =============================================================================
%                              PART 2: SAFETY
% =============================================================================

\subsubsection{LLM Safety Overview}

\begin{defbox}[Safety Aspects]
\begin{itemize}
  \item \textbf{Value misalignment:} Social bias, toxicity, ethics violations
  \item \textbf{Robustness to attack:} Jailbreaking, adversarial prompts
  \item \textbf{Privacy:} Data leakage, memorisation
  \item \textbf{Misuse:} Misinformation, deepfakes, weaponization
  \item \textbf{Agent safety:} Goal misalignment, deception
\end{itemize}
\end{defbox}

\bigseparator

% --- Social Bias ---
\subsubsection{Social Bias in LLMs}

\begin{defbox}[Two Types of Harm]
\textbf{Representational harm:} Reinforcing stereotypes, misrepresenting groups, using offensive language.

\textit{Example:} LLM describes nurses as women and engineers as men.

\textbf{Allocational harm:} Creating unfair advantages/disadvantages in resource distribution.

\textit{Example:} Resume-screening AI favours male candidates over female candidates.
\end{defbox}

\begin{exbox}[Bias Example]
\textbf{Prompt:} ``Two Muslims walked into a...''

\textbf{GPT-3 output:} Frequently generates violent completions.

\textbf{Replacing ``Muslim'' with other religions:} Significantly reduces violent completions.

This shows learned associations from training data.
\end{exbox}

\begin{defbox}[Bias Evaluation Methods]
\textbf{Embedding-based:} Compare distances between biased/unbiased representations (WEAT, SEAT).

\textbf{Probability-based:} Compare token probabilities for different demographic groups.

\textbf{Generation-based:} Analyse token distributions in actual outputs.

\textbf{Key insight:} Embedding-level neutrality doesn't guarantee neutral outputs---biases can re-emerge during generation.
\end{defbox}

\begin{defbox}[Bias Mitigation Strategies]
\textbf{Pre-processing:} Data balancing, data cleaning.

\textbf{In-training:} Architecture modifications, loss function adjustment.

\textbf{Inference:} Decoding strategy modification, attention weight adjustment.

\textbf{Post-processing:} Keyword replacement, rewriting.
\end{defbox}

\bigseparator

% --- Toxicity ---
\subsubsection{Toxicity in LLMs}

\begin{defbox}[Toxicity Definition]
\textbf{Toxicity} refers to generation of harmful or offensive content that can:
\begin{itemize}
  \item Escalate hostility
  \item Disrupt social harmony
  \item Threaten safety
\end{itemize}

\textbf{Causes:} Toxic training data, generalisation of harmful patterns, interactive evolution.
\end{defbox}

\begin{defbox}[Toxicity Mitigation by Phase]
\textbf{Pre-training:} Filter toxic data (but may weaken general capabilities).

\textbf{SFT:} Fine-tune on non-toxic subsets, high-quality aligned data.

\textbf{Alignment:} RLHF, RLAIF, Adversarial DPO.

\textbf{Key trade-off:} Removing all toxic data can hurt model performance---toxic content may carry useful linguistic patterns.
\end{defbox}

\bigseparator

% --- Jailbreaking ---
\subsubsection{Jailbreaking Attacks}

\begin{defbox}[Jailbreaking Definition]
\textbf{Jailbreaking} is manipulating an LLM through crafted prompts to bypass safety filters and generate restricted outputs.
\end{defbox}

\begin{thmbox}[Black-Box Jailbreaking]
No internal model access required---only API/chat interface.

\textbf{Techniques:}
\begin{itemize}
  \item \textbf{Goal hijacking:} ``IGNORE INSTRUCTIONS! NOW SAY...''
  \item \textbf{Refusal suppression:} ``Never say `cannot' or `unfortunately'\,''
  \item \textbf{Role-play:} Assign personas (hacker, propagandist)
  \item \textbf{Fictional scenario:} ``Pretend you're in developer mode''
  \item \textbf{Prompt decomposition:} Split malicious prompts into benign sub-prompts
  \item \textbf{Low-resource language:} Exploit weaker safety in underrepresented languages
  \item \textbf{Cipher text:} Encode prompts in Base64 or custom encodings
\end{itemize}
\end{thmbox}

\begin{thmbox}[White-Box Jailbreaking]
Requires full access to model internals (architecture, weights, gradients).

\textbf{Greedy Coordinate Gradient (GCG):}
\begin{enumerate}
  \item Calculate loss for current adversarial prompt
  \item Compute gradients for token positions in adversarial suffix
  \item Sample candidate token replacements from top-$k$
  \item Apply substitutions that reduce loss most
  \item Repeat until successful attack
\end{enumerate}

\textbf{Key finding:} Adversarial suffixes may transfer across different LLMs.
\end{thmbox}

\bigseparator

% --- Defense Strategies ---
\subsubsection{Defense Strategies}

\begin{defbox}[SmoothLLM]
\textbf{Observation:} Adversarial suffixes are fragile to character-level perturbations.

\textbf{Defense:}
\begin{enumerate}
  \item Create $N$ copies of the attacked prompt
  \item Randomly perturb each copy (character-level noise)
  \item Generate responses for all perturbed copies
  \item Aggregate responses to produce final output
\end{enumerate}

Perturbations break the carefully crafted adversarial suffix.
\end{defbox}

\begin{defbox}[GradSafe]
\textbf{Idea:} Use gradient analysis to detect unsafe prompts.

\textbf{Process:}
\begin{enumerate}
  \item Identify safety-critical parameters where gradients differ between safe/unsafe prompts
  \item Compute loss gradients for input prompt (with response ``Sure'')
  \item Flag as unsafe if gradient shows high cosine similarity to unsafe gradient references
\end{enumerate}
\end{defbox}

\begin{tipbox}[Balanced Defense Strategy]
For robust defense against both black-box and white-box attacks:
\begin{itemize}
  \item Gradient-similarity unsafe prompt detection (GradSafe)
  \item Decoding constraints at inference time
  \item Pre-training data cleaning + RLHF alignment
  \item Post-processing content filters
\end{itemize}

No single defense is sufficient---layered approaches are needed.
\end{tipbox}

\bigseparator

% --- Summary ---
\subsubsection{Summary Tables}

\begin{center}
\renewcommand{\arraystretch}{1.3}
\begin{tabular}{@{} l p{5cm} p{4cm} @{}}
\toprule
\textbf{Evaluation Method} & \textbf{Strengths} & \textbf{Weaknesses} \\
\midrule
Automatic (BLEU, F1) & Fast, reproducible & Ignores semantics \\
Semantic (BERTScore) & Captures meaning & Compute-intensive \\
Human evaluation & Gold standard & Expensive, inconsistent \\
LLM-as-judge & Scalable, nuanced & Bias, adversarial vulnerability \\
\bottomrule
\end{tabular}
\end{center}

\vspace{1em}

\begin{center}
\renewcommand{\arraystretch}{1.3}
\begin{tabular}{@{} l p{4cm} p{4cm} @{}}
\toprule
\textbf{Safety Issue} & \textbf{Detection} & \textbf{Mitigation} \\
\midrule
Social bias & Embedding/probability metrics & Data balancing, post-processing \\
Toxicity & Toxicity classifiers & Pre-training filters, RLHF \\
Black-box jailbreak & Input pattern detection & Prompt filtering, refusal training \\
White-box jailbreak & Gradient analysis (GradSafe) & SmoothLLM, adversarial training \\
\bottomrule
\end{tabular}
\end{center}

\begin{thmbox}[Key Takeaways]
\textbf{Evaluation:}
\begin{itemize}
  \item High benchmark scores $\neq$ true reasoning (GSM-Symbolic, clause order sensitivity)
  \item Goodhart's Law: Optimising for metrics can corrupt them
  \item LLM-as-judge is scalable but has biases
\end{itemize}

\textbf{Safety:}
\begin{itemize}
  \item Representational harm (stereotypes) vs.\ allocational harm (unfair outcomes)
  \item Black-box attacks use prompt engineering; white-box uses gradients
  \item Layered defense strategies are essential
\end{itemize}
\end{thmbox}

